% Generated draft: requires mathematical and source review.
\documentclass[11pt]{article}
\usepackage[margin=25mm]{geometry}
\usepackage{fontspec}
\setmainfont{Latin Modern Roman}
\usepackage{amsmath,amssymb,mathtools,unicode-math}
\setmathfont{Latin Modern Math}
\usepackage{enumitem,needspace,xurl,hyperref}
\setlength{\parindent}{0pt}
\setlength{\parskip}{5pt}
\setlength{\emergencystretch}{3em}
\newcommand{\N}{\mathbb{N}}
\newcommand{\Z}{\mathbb{Z}}
\newcommand{\Q}{\mathbb{Q}}
\newcommand{\R}{\mathbb{R}}
\newcommand{\C}{\mathbb{C}}
\newcommand{\F}{\mathbb{F}}
\newcommand{\A}{\mathbb{A}}
\newcommand{\PP}{\mathbb{P}}
\newcommand{\E}{\mathbb{E}}
\newcommand{\eps}{\varepsilon}
\newcommand{\dd}{\,\mathrm{d}}
\begin{document}
\section*{UnsolvedMath ID: 30001461}
\textbf{Source problem code:} catalogue entry
\subsection{Definitions and mathematical statement}
Write \(U(\mathfrak g,e)\) for the finite \(W\)-algebra associated with a Lie algebra \(\mathfrak g\) and a nilpotent element \(e\in\mathfrak g\), once the ground field, admissible hypotheses, and construction have been specified. This notation identifies the usual input data; the supplied record does not specify those choices.

\textbf{Source problem.} Determine the structure and representation theory of finite \(W\)-algebras associated with classical Lie algebras.

\textbf{Unresolved scope.} The record does not specify the ground field or its characteristic, the precise class of classical Lie algebras, the nilpotent elements under consideration, or the intended categories of representations. It also does not specify what constitutes determining the structure or representation theory, such as presentations, classification of modules, or character formulas. Consequently, this is a research programme rather than a uniquely specified proposition or classification problem. No additional hypotheses, quantifiers, or completion criterion are imposed in this draft.
\subsection{Short English statement}
The problem asks for an understanding of the algebras and representations of finite \(W\)-algebras attached to classical Lie algebras. Its intended scope and required outputs remain unspecified, so a precise classification target requires human clarification.
\subsection{Sources}
\begin{itemize}
\item Local UnsolvedMath registry, permanent ID \(30001461\), title \emph{Representation Theory of Classical Finite W-Algebras}, source code \texttt{catalogue entry}, in the Oberwolfach Reports Open Problems collection. The record gives the proposed year as 2010 and supplies no individual proposer.
\item \emph{Algebraic Groups} (2010), Oberwolfach Reports, DOI 10.4171/owr/2010/19, pp. 1101--1163: \url{https://doi.org/10.4171/owr/2010/19}. This is also the supplied literature source labelled ``original OWR report.'' The background describes Oberwolfach Reports as documenting mathematical workshops and their open-problem discussions.
\item Supplied literature source labelled ``Futorny--Molev--Ovsienko'': \url{https://arxiv.org/abs/0801.0903}. The record supplies no specific theorem or passage connecting this reference to a precise version of the problem.
\item The registry labels the statement ``exact'' and reports that its mathematical content was checked against the official Oberwolfach report, with notation, context, and question form normalized. These are supplied verification claims, not independent verification performed for this draft.
\item Literature assessment dated \texttt{2026-08-21}: the supplied curation labels the problem open and describes it as a broad programme, with later results addressing substantial subclasses and representation-theoretic questions but not the general structure and representation theory. This is supplied dated curation, not an independently verified literature assessment; its time-sensitive status and cited claims require independent verification.
\end{itemize}
\end{document}
